\section{End-Term Comprehensive Examination --- Set A}
\label{sec:endterm_set_a}

\LPUPaperHeader{End-Term Comprehensive Examination}{A}{180 Minutes (3 Hours)}{50 Marks}{CO1--CO6 (Comprehensive Syllabus Units 1 to 6)}

\subsection*{Section A: Objective Multiple Choice Questions [10 Marks]}
\begin{center}
\sffamily\textbf{\color{AlertRed}Attempt all ten questions. Each question carries 1 mark.}
\end{center}

\mcqitem{1}{Which problem solving approach is best suited for complex, non-linear human problems?}{Design Thinking}{Linear Waterfall Method}{Trial-and-Error without reflection}{Premature Optimization}{1}

\mcqitem{2}{In the Stanford d.school framework, what is the primary deliverable of the Define stage?}{A functional codebase}{Point-of-View (POV) statement}{An invoice for services}{A usability survey}{1}

\mcqitem{3}{Which tool maps user emotions across touchpoints over time?}{Fishbone Diagram}{Customer Journey Map}{Pareto Chart}{Gantt Chart}{1}

\mcqitem{4}{What does the 'A' in the SCAMPER creative ideation methodology stand for?}{Analyze}{Adapt}{Automate}{Accelerate}{1}

\mcqitem{5}{In Brainwriting 6-3-5, how many total ideas are produced after 6 rounds?}{36}{72}{108}{216}{1}

\mcqitem{6}{What is the primary characteristic of a Wizard of Oz prototype?}{It uses expensive robotics}{The front-end appears automated while backend tasks are done manually}{It only exists as code without UI}{It is tested without real users}{1}

\mcqitem{7}{An industry average System Usability Scale (SUS) score is approximately:}{30.0}{50.0}{68.0}{95.0}{1}

\mcqitem{8}{In Barbara Minto's SCQA presentation framework, what does 'C' stand for?}{Conclusion}{Complication}{Context}{Criteria}{1}

\mcqitem{9}{What is the first stage in David Kolb's Experiential Learning Cycle?}{Concrete Experience}{Reflective Observation}{Abstract Conceptualization}{Active Experimentation}{1}

\mcqitem{10}{In a blameless postmortem, errors are primarily attributed to:}{Careless junior employees}{Systemic and architectural flaws}{Malicious intent}{Bad luck}{1}

\newpage
\subsection*{Section B: Applied Analytical Engineering Scenarios [20 Marks]}
\begin{center}
\sffamily\textbf{\color{AlertRed}Attempt all four questions. Each question carries 5 marks.}
\end{center}

\questionitem{11}{
    A smart mobility startup is selecting an electric battery chemistry for electric fleet delivery scooters. The criteria and weights are: Energy Density ($w_1 = 0.35$), Thermal Safety ($w_2 = 0.30$), Lifecycle Degradation ($w_3 = 0.20$), and Sourcing Cost ($w_4 = 0.15$). The ratings (1--10) are:
    \begin{itemize}
        \item \textbf{Chemistry A (NMC):} Density = 9, Safety = 6, Lifecycle = 7, Cost = 6
        \item \textbf{Chemistry B (LFP):} Density = 7, Safety = 9, Lifecycle = 9, Cost = 8
    \end{itemize}
    Execute an MCDA calculation to determine the optimal battery chemistry for the commercial fleet.
}{5}{CO3}{L3 Apply}

\questionitem{12}{
    Deconstruct the Root Causes of high student drop-out rates in online asynchronous computer programming courses using an Ishikawa (Fishbone) diagram across the 4Ps (People, Policy, Procedures, Plant/Technology).
}{5}{CO2}{L4 Analyze}

\questionitem{13}{
    Construct a Multi-Tier Service Blueprint for a self-service electric vehicle battery swapping station. Include Customer Actions, Frontstage, Backstage, Support Processes, and at least two Failure Nodes [!].
}{5}{CO5}{L4 Analyze}

\questionitem{14}{
    Structure an Executive Pitch Presentation using Barbara Minto's Pyramid Principle (SCQA) proposing to replace passwords with FIDO2 hardware passkeys in an enterprise bank.
}{5}{CO6}{L3 Apply}

\newpage
\subsection*{Section C: Comprehensive Capstone Design Challenges [20 Marks]}
\begin{center}
\sffamily\textbf{\color{AlertRed}Attempt both comprehensive design challenges. Each question carries 10 marks.}
\end{center}

\questionitem{15}{
    \textbf{Design Challenge 1: Urban Public Healthcare Access}\\
    An overburdened regional public health department experiences severe outpatient clinic congestion, long queues, and widespread medication shortages. Elderly patients often travel three hours by bus only to discover that prescribed hypertension medications are out of stock.\\
    \textbf{Deliverables Required:}\\
    (a) Analyze why this scenario qualifies as an ill-structured wicked problem (3 Marks).\\
    (b) Formulate a human-centered Point-of-View (POV) statement and derive three actionable HMW questions (3 Marks).\\
    (c) Apply the SCAMPER method to generate four innovative non-capital-intensive interventions (4 Marks).
}{10}{CO1, CO2, CO4}{L5 Evaluate}

\questionitem{16}{
    \textbf{Design Challenge 2: Post-Incident Reflective Audit \& Validation Architecture}\\
    A fintech peer-to-peer payments app launched a major redesign that caused a 40\% drop in transaction volume and an influx of 15,000 angry customer support tickets complaining that the ``Send Money'' button had disappeared into an obscure nested menu.\\
    \textbf{Deliverables Required:}\\
    (a) Conduct an incident reflective audit using Graham Gibbs' Reflective Cycle across all six stages (5 Marks).\\
    (b) Design a rapid 14-day usability validation sprint utilizing the System Usability Scale (SUS) and a 4-Quadrant Feedback Capture Grid to rectify the interface (5 Marks).
}{10}{CO5, CO6}{L5 Evaluate}

\newpage
\subsection*{Solutions and Marking Rubrics --- End-Term Set A}

\begin{solutionbox}[Solutions to Section A: Objective MCQs]
\textbf{1.} (A) Design Thinking \quad
\textbf{2.} (B) Point-of-View (POV) statement \quad
\textbf{3.} (B) Customer Journey Map \quad
\textbf{4.} (B) Adapt \quad
\textbf{5.} (C) 108 \quad
\textbf{6.} (B) Front-end appears automated while backend tasks are done manually \quad
\textbf{7.} (C) 68.0 \quad
\textbf{8.} (B) Complication \quad
\textbf{9.} (A) Concrete Experience \quad
\textbf{10.} (B) Systemic and architectural flaws
\begin{markingrubric}
1 Mark per correct MCQ = \textbf{10 Marks Total}.
\end{markingrubric}
\end{solutionbox}

\begin{solutionbox}[Solution to Question 11: Battery Chemistry MCDA]
Calculate $S_i = \sum w_j \cdot r_{ij}$:\\
$\bullet$ \textbf{Chemistry A (NMC):}\\
$S_A = (0.35 \times 9) + (0.30 \times 6) + (0.20 \times 7) + (0.15 \times 6) = 3.15 + 1.80 + 1.40 + 0.90 = \mathbf{7.25}$\\
$\bullet$ \textbf{Chemistry B (LFP):}\\
$S_B = (0.35 \times 7) + (0.30 \times 9) + (0.20 \times 9) + (0.15 \times 8) = 2.45 + 2.70 + 1.80 + 1.20 = \mathbf{8.15}$\\
\textbf{Selection Conclusion:}\\
Chemistry B (LFP) decisively outperforms Chemistry A ($8.15$ vs $7.25$), because commercial delivery fleets prioritize thermal safety, long cycle life, and low unit costs over pure energy density.
\begin{markingrubric}
$\bullet$ Accurate calculation for NMC and LFP: \textbf{3 Marks}\\
$\bullet$ Clear analytical justification and final selection: \textbf{2 Marks}
\end{markingrubric}
\end{solutionbox}

\begin{solutionbox}[Solution to Question 12: Ishikawa 4P for Coding Course Drop-out]
$\bullet$ \textbf{People:} Imposter syndrome among beginners; cognitive overload from syntax; lack of peer belonging; unresponsive teaching assistants.\\
$\bullet$ \textbf{Policy:} Rigid non-negotiable weekly assignment deadlines; all-or-nothing grading; zero grace periods for working adults.\\
$\bullet$ \textbf{Procedures:} Passive 90-minute lecture videos without active coding checkpoints; lack of automated scaffolding hints.\\
$\bullet$ \textbf{Plant / Technology:} Clunky web IDE environment that fails to save code states; confusing error compiler messages; lack of dark mode or mobile editor support.
\begin{markingrubric}
$\bullet$ Inclusion of all 4P categories: \textbf{2 Marks}\\
$\bullet$ Deep, realistic pedagogical and technical insights: \textbf{3 Marks}
\end{markingrubric}
\end{solutionbox}

\begin{solutionbox}[Solution to Question 13: Battery Swapping Service Blueprint]
$\bullet$ \textbf{Physical Evidence:} Kiosk touchscreen, automated mechanical compartment door, status LED, invoice notification.\\
$\bullet$ \textbf{Customer Actions:} Pulls scooter into bay $\to$ scans QR code $\to$ removes depleted battery $\to$ inserts into slot $\to$ retrieves fresh battery $\to$ clicks into scooter.\\
\textbf{--- Line of Interaction ---}\\
$\bullet$ \textbf{Frontstage Contact:} Automated voice guidance from kiosk; mechanical locking mechanism releasing charged battery.\\
\textbf{--- Line of Visibility ---}\\
$\bullet$ \textbf{Backstage Contact:} Internal conveyor belts moving battery to high-voltage cooling rack; robotic state-of-health (SoH) diagnostics.\\
\textbf{--- Line of Internal Interaction ---}\\
$\bullet$ \textbf{Support Processes:} Cloud billing API deduction; grid load balancer managing off-peak thermal charging.\\
$\bullet$ \textbf{Failure Nodes [!]:}\\
1. \textit{[! 1] Lock Jam:} Battery compartment door fails to unlatch due to thermal expansion.\\
2. \textit{[! 2] Swapped Battery Degradation:} Cell with degraded capacity ($\le 60\%$ SoH) inadvertently released to rider.
\begin{markingrubric}
$\bullet$ Complete blueprint architecture across interaction and visibility lines: \textbf{3 Marks}\\
$\bullet$ Identification and mitigation of 2 critical failure nodes: \textbf{2 Marks}
\end{markingrubric}
\end{solutionbox}

\begin{solutionbox}[Solution to Question 14: Minto SCQA Passkeys Executive Pitch]
$\bullet$ \textbf{Situation (S):} Enterprise bank serves 4 million active mobile users handling \$2.8B in annual transactions using SMS OTP and complex passwords.\\
$\bullet$ \textbf{Complication (C):} Phishing attacks and SIM-swapping fraud increased by 140\% this year, inflicting \$18M in customer fraud reimbursements and generating severe reputational damage.\\
$\bullet$ \textbf{Question (Q):} How can we eliminate credential phishing and reduce login friction without increasing customer drop-off?\\
$\bullet$ \textbf{Answer (A):} Deploy biometric FIDO2 hardware passkeys natively integrated with FaceID/fingerprint biometrics, reducing login latency by 70\% and mathematically eliminating credential interception.
\begin{markingrubric}
$\bullet$ Complete SCQA formulation with business context: \textbf{3 Marks}\\
$\bullet$ Persuasive executive alignment connecting security to user experience: \textbf{2 Marks}
\end{markingrubric}
\end{solutionbox}

\begin{solutionbox}[Solution to Question 15: Public Healthcare Capstone Challenge]
\textbf{(a) Wicked Problem Analysis (3 Marks):}\\
1. \textit{No definitive formulation:} The challenge intertwines medical supply chains, municipal bus schedules, rural poverty, and pharmaceutical inventory policies.\\
2. \textit{Irreversible consequences:} Stock-outs cause stroke, cardiac arrest, and death; interventions cannot be tested trivially in sandboxes.\\
3. \textit{Competing stakeholders:} Doctors, bus operators, pharmaceutical distributors, and patients hold conflicting objectives.\\
\textbf{(b) POV Statement \& HMW Questions (3 Marks):}\\
$\bullet$ \textit{POV:} ``Elderly hypertensive patients [User] need guaranteed medication availability and decentralized local dispensing [Need] because exhausting transit journeys ending in empty pharmacy shelves cause life-threatening health crises and loss of trust [Insight].''\\
$\bullet$ \textit{HMW 1:} How might we communicate real-time pharmacy medicine stock via basic SMS before patients leave home?\\
$\bullet$ \textit{HMW 2:} How might we utilize existing rural post offices as secure monthly medication pick-up points?\\
$\bullet$ \textit{HMW 3:} How might we deliver 90-day supply batches to chronic patients to reduce travel frequency by 66\%?\\
\textbf{(c) SCAMPER Ideation Concepts (4 Marks):}\\
$\bullet$ \textbf{Substitute (S):} Substitute in-person visits with automated vernacular phone calls for routine prescription refills.\\
$\bullet$ \textbf{Combine (C):} Combine public postal delivery vans with pharmaceutical cold-chain parcel drops.\\
$\bullet$ \textbf{Adapt (A):} Adapt the Amazon hub locker model by placing smart medication lockers at rural bus junctions.\\
$\bullet$ \textbf{Eliminate (E):} Eliminate paper prescription renewals by auto-authorizing chronic hypertension refills upon verified blood pressure readings.
\begin{markingrubric}
$\bullet$ Part (a) Wicked problem framing: \textbf{3 Marks}\\
$\bullet$ Part (b) POV statement and 3 HMW questions: \textbf{3 Marks}\\
$\bullet$ Part (c) 4 SCAMPER interventions: \textbf{4 Marks}
\end{markingrubric}
\end{solutionbox}

\begin{solutionbox}[Solution to Question 16: Gibbs Cycle Audit \& Usability Sprint]
\textbf{(a) Gibbs' Reflective Cycle Deployment (5 Marks):}\\
1. \textbf{Description:} App redesign launched; transaction volume crashed 40\%; 15k tickets received complaining that 'Send Money' was buried.\\
2. \textbf{Feelings:} Team felt defensive, arguing that the redesign was 'aesthetically cleaner'; then panic set in as revenue dropped.\\
3. \textbf{Evaluation:} Aesthetics improved, but core usability and discoverability were catastrophic failures.\\
4. \textbf{Analysis:} The team fell victim to \textit{Designer Vanity} and confirmation bias, testing the UI only with internal designers who already knew where buttons were hidden. Zero think-aloud usability testing was conducted with authentic users.\\
5. \textbf{Conclusion:} Clean aesthetics must never compromise task affordances. Core transactional triggers must remain primary on home screens.\\
6. \textbf{Action Plan:} Immediately rollback navigation; implement mandatory usability gating ($\text{SUS} \ge 80$) before any production release.\\
\textbf{(b) 14-Day Validation Sprint Architecture (5 Marks):}\\
$\bullet$ \textit{Days 1--3:} Rapid wireframe refactor restoring prominent 'Send Money' action button on bottom navigation dock.\\
$\bullet$ \textit{Days 4--7:} Conduct moderated usability testing with 15 representative users, tracking Task Completion Rate (TCR) and computing composite System Usability Scale (SUS) scores.\\
$\bullet$ \textit{Days 8--10:} Populate Feedback Capture Grid: synthesize remaining Likes, Criticisms, and Questions.\\
$\bullet$ \textit{Days 11--14:} Verify that $\text{SUS} \ge 82.0$ and $\text{TCR} = 100\%$; deploy phased 10\% canary release to production.
\begin{markingrubric}
$\bullet$ Part (a) All 6 stages of Gibbs Cycle rigorously articulated: \textbf{5 Marks}\\
$\bullet$ Part (b) 14-day sprint protocol with SUS metrics and feedback grid: \textbf{5 Marks}
\end{markingrubric}
\end{solutionbox}
